\documentclass[10pt]{article}
\usepackage{resume}
\usepackage{xeCJK}
\setCJKmainfont{FandolSong-Regular.otf}[BoldFont=FandolSong-Bold.otf,Script=Default]
\hypersetup{pdftitle={李哲骧 - 研究岗位简历},pdflang={zh-CN}}
\begin{document}
\chineseheader

\section{教育背景}
\entry{加州大学洛杉矶分校（UCLA）}{2026.09 - 2028.06（预计）}
计算机科学硕士在读
\entry{南加州大学（USC）}{2026.05 毕业}
计算机科学学士；GPA：3.76/4.00；Viterbi Dean's List

\section{研究技能}
\textbf{研究方法：}多模态学习、高效注意力、视觉 token 筛选、分子性质预测\\
\textbf{实验工具：}Python、PyTorch、Triton、RDKit、Chemprop、C/C++、SQL、Linux\\
\textbf{评估方法：}患者级数据划分、分子骨架划分、难负样本构建、消融实验、SHAP 分析

\section{科研经历}
\entry{UCLA Health Intelligence Lab（HAIL）｜研究生研究员}{2026.08 - 至今}
\detail{导师：Yuzhe Yang 教授}
\begin{itemize}
\item 围绕糖尿病血糖预测开展多模态基础模型研究，将连续血糖监测（CGM）与胰岛素、饮食、活动及患者背景信息结合。
\item 搭建 MetaboNet 基准与数据检查流程，按患者划分数据；比较 GluFormer、T1DFormer、自研预测状态模型和 Transformer 在未来 30、60、120 分钟的预测表现。
\end{itemize}
\entry{USC Yan Liu 教授课题组｜药物发现研究}{2026 年春}
\begin{itemize}
\item 跟随课题组博士生开展抗病毒药物发现研究，在七个 TDC 数据集上比较分子描述符模型、Chemprop 与混合模型，采用分子骨架划分检验泛化能力。
\item 分析 RDKit 描述符重要性和 SHAP 归因，考察模型对抗病毒活性的预测主要依据哪些分子特征。
\end{itemize}

\section{代表论文与研究}
\entry{MedTokenBudget：保留病灶信息的视觉 token 路由}{2026}
\detail{Zhexiang Li；预印本。\link{https://arxiv.org/abs/2609.30613}{arXiv:2609.30613}}
\begin{itemize}
\item 在视觉骨干网络后加入病灶监督路由器，结合预算课程学习、多样性约束和注意力蒸馏训练。
\item 在 ISIC 2019 上统计固定 token 预算下的标注病灶保留比例，检查筛选过程是否丢失关键区域。
\end{itemize}
\entry{基于质量守恒证据的 MS/MS 谱图检索}{2026}
\detail{稿件已投稿 AAAI 2027}
\begin{itemize}
\item 将满足质量守恒关系的碎片对编码为稀疏 token；构建质量接近和谱峰重叠的难负样本，并通过证据移除实验检验检索排序的依据。
\end{itemize}
\entry{Efficient Attention via Pre-Scoring：注意力预评分}{2025}
\detail{第一作者；与 H. Wang、Y. Bao、D. Woodruff 合作。审稿中。\link{https://arxiv.org/abs/2505.11040}{arXiv:2505.11040}}
\begin{itemize}
\item 提出先对 key 区间进行预评分、再计算注意力的方法，并实现 Triton 内核。
\item 与 HyperAttention、LevAttention 对比评估，考察不同注意力近似方法的表现。
\end{itemize}

\section{相关实践}
\entry{未来求索 / PMAOS｜AI 工程实习生}{2026.07 - 2026.08}
\begin{itemize}
\item 设计会话、桥接和长期三档推荐曝光预算；开发支持 11 种语言的音视频情绪识别，在模态证据冲突时暂不判断。
\item 为算法实现建立验证依据：推荐部分完成 629 项测试，多模态部分通过 139 项 Python 测试。
\end{itemize}
\entry{UCSC 机器学习导论｜课程辅导员}{2023 年春}
\begin{itemize}
\item 在 Lise Getoor 教授课程中讲解逻辑回归、SVM 和树模型，共同设计 Jupyter 作业，并参与编程题自动评分器开发。
\end{itemize}
\end{document}
